\chapter*{Acknowledgments}
\linespread{1.12}\selectfont

A la profesora {\bfseries\color{darkred}Raquel Pezoa}, quien me ha guiado y
ayudado con absolutamente todo, desde mi licenciatura en física hasta esta tesis.
Nada de esta investigación existiría sin su acompañamiento; su disposición es
admirable y la llevaré siempre conmigo.

Al profesor {\bfseries\color{darkred}Nicolás Viaux}, por guiarme desde el lado
más teórico con ideas y consejos que enriquecieron este trabajo.

A {\bfseries\color{darkred}Liza}, mi gran compañera de vida, que siempre me ha
apoyado, comprendido y levantado, sin importar la situación. Es mi inspiración
para ser mejor persona, mejor investigador y alguien completamente feliz. La
admiro y la amo desde lo más profundo de mi corazón.

A {\bfseries\color{darkred}Yerko} y {\bfseries\color{darkred}Daniel},
compañeros de física desde los inicios y mis dos pokémon starters, por su
sabiduría y compañerismo durante toda esta travesía. A mis amigos de toda la
vida ---{\bfseries\color{darkred}Naxo}, {\bfseries\color{darkred}Nico},
{\bfseries\color{darkred}Julito}, {\bfseries\color{darkred}Tatis},
{\bfseries\color{darkred}Maury} y {\bfseries\color{darkred}Claudia}---, que
forjaron quién soy y llenan de alegría mi vida.

A mis padres, {\bfseries\color{darkred}Marcelo} y
{\bfseries\color{darkred}Ximena}, a quienes amo con todo mi corazón y estaré
agradecido de por vida: son el pilar que me permitió llegar hasta donde estoy y
ser quien soy. A mis hermanos, {\bfseries\color{darkred}Nicolás},
{\bfseries\color{darkred}Marcelo} y {\bfseries\color{darkred}Jesús}, personas
inteligentísimas que, cada uno a su manera, han nutrido mi conocimiento
profesional y de la vida.

Finalmente, a mi gatito {\bfseries\color{darkred}Thorfinn}, que sin hablar me
consuela, con un maullido alegra mi día entero y demuestra su amor con rasguños
y ganas de comer. Me espera cada día e hizo que mis días nunca más fuesen
solitarios; me demostró que no se necesitan palabras para querer.

% Separator
\begin{center}
\begin{tikzpicture}[baseline=-0.5ex]
  \draw[darkred, line width=0.4pt] (-3.2,0) -- (-0.22,0);
  \draw[darkred, line width=0.4pt] (0.22,0) -- (3.2,0);
  \fill[darkred] (0,0.09) -- (0.09,0) -- (0,-0.09) -- (-0.09,0) -- cycle;
  \fill[darkred] (-0.17,0) circle (0.02);
  \fill[darkred] (0.17,0) circle (0.02);
\end{tikzpicture}
\end{center}

{\linespread{1.0}\selectfont\small This work was supported by ANID CCTVal (CIA250027) and by Proyectos Internos de Investigación USM 2025 (PI\_LII\_25\_06). The author also acknowledges the support of the Becas Magíster Científico-Tecnológico of the Universidad Técnica Federico Santa María.\par}
